\section{阅读地图：从 EP137 到后续队列}

EP137 在Zhang Xiaojun AI 队列里的位置很特殊。EP140 和 EP138 更偏模型实验室与 Agent 产业范式；EP139 提供 Agent 技术史；EP137 则把问题推向科学研究与形式化证明。它是从“AI 改变生产力”到“AI 改变知识生产”的桥。

\subsection{读者应该带走的五个问题}

第一，AI 能否从题目求解走向真正的数学研究？第二，Lean 是否会成为 AI for Math 的主接口？第三，数学家直觉能否被模型学习，还是必须长期由人类提供？第四，Axiom 这类 Neo Lab 如何在基础研究和商业化之间找到中间产物？第五，AI for Math 的成功是否会反过来推动 AI 自身算法和理论进步？

\begin{knowledgebox}{本期与后续整理的关系}
后续如果整理更多 AI for Science、机器人、世界模型或自动驾驶访谈，可以复用本期框架：先找可验证环境，再看人类专家反馈，再看模型是否能形成长期闭环。数学只是这个框架中验证最硬的一类场景。
\end{knowledgebox}

\subsection{本章小结}

EP137 的价值在于拓展了 AI 队列的边界：从 Agent 产品和模型训练，进入数学和科学发现。它提醒我们，最难的 AI 应用往往同时需要形式验证和人类审美。


